\documentclass[11pt,a4paper]{article}
\usepackage[margin=2cm]{geometry}
\usepackage{fontspec}
\usepackage{xeCJK}
\setCJKmainfont{Songti TC}
\setCJKsansfont{PingFang TC}
\usepackage{enumitem,xcolor,parskip}
\usepackage{amssymb}
\setlist{nosep,leftmargin=1.2em}
\newcommand{\tbd}[1]{\textcolor{red}{[#1]}}
\pagestyle{empty}
\begin{document}
\begin{center}{\Large\bfseries Information and Consent Form}\\[2pt]
SeeSound pronunciation recording study\end{center}

\textbf{Who we are.} A research team at Ho Chi Minh City University of Foreign Languages - Information Technology (HUFLIT) and National Formosa University. Contact: Prof.\ Hui-Kai Su, hksu@nfu.edu.tw. The data are managed by Kiet Pham (HUFLIT).

\textbf{Why we ask.} We are testing a web program (SeeSound) that shows deaf and hard-of-hearing learners how to say English consonants. Before it is used with children, we want to check, with teachers, how well the program's judgements of a spoken word agree with what teachers hear. This is a small pilot.

\textbf{What happens.} You (or your child) say eight English words (\textit{she, see, fee, pea, bee, me, knee, cheat}) twice each into a computer microphone. This takes about 4 minutes. The screen does not show a score. A teacher may later listen to the recordings and rate them.

\textbf{What we keep.} A participant code (not a name), the recordings \emph{only if you tick the box below}, and numbers measured from them. No camera or video is used. Data stay on a research computer managed by Kiet Pham; recordings are heard only by the raters and the research team, are not sent by e-mail or chat, and are deleted after 1 year. Reports show combined results only.

\textbf{Risks and benefits.} No known risks beyond a few minutes of speaking. There is no payment and no direct benefit; the study may help improve tools for learners.

\textbf{Your rights.} Taking part is voluntary. You or your child may stop at any time, and you may ask us to delete the recordings and data, without giving a reason and without any effect on your teaching or care. To withdraw, give the participant code and contact us at the address above.

\medskip\hrule\medskip
\textbf{Consent} (for a child, the parent or guardian completes this; the child is also asked whether they are happy to take part)

\begin{itemize}[label=$\square$]
\item I have read (or had read to me) this information and could ask questions.
\item I agree that I / my child takes part in the recording described above.
\item I \textbf{also} agree that the \textbf{audio recordings} are kept and listened to by the raters and research team as described. \textit{(If this is not ticked, the person can still take part, but teachers cannot rate the recordings.)}
\end{itemize}

\vspace{6pt}
\newcommand{\line}[2]{\noindent #1 \hrulefill\hspace{1em}#2 \rule{4cm}{0.4pt}\\[12pt]}
\noindent Participant code (filled in by the teacher): \rule{3.5cm}{0.4pt}\hfill Date: \rule{4cm}{0.4pt}\\[12pt]
\noindent Name of adult participant or parent/guardian: \hrulefill\\[12pt]
\noindent Signature: \hrulefill\\[12pt]
\noindent Teacher's initials (who explained the study): \rule{3cm}{0.4pt}\\[12pt]
\noindent Child's assent (``I am happy to do this''), if able: \rule{4cm}{0.4pt}

\vspace{4pt}\footnotesize Keep this form separate from the recordings: the form has a name, the data have only the code.

\newpage
\begin{center}{\Large\bfseries 研究說明暨同意書}\\[2pt]
SeeSound 發音錄音研究\end{center}

\textbf{研究團隊。} 本研究由胡志明市外語暨資訊大學（HUFLIT）與國立虎尾科技大學的研究團隊進行。聯絡人：蘇暉凱教授，hksu@nfu.edu.tw。資料由 Kiet Pham（HUFLIT）負責管理。

\textbf{為什麼邀請您。} 我們正在測試一個網頁程式（SeeSound），用來幫助聽障與重聽的學習者看見並練習英語子音。在這個程式用於兒童之前，我們想請老師協助確認：程式對一個單字發音的判斷，與老師聽到的結果有多一致。這是一個小規模的前導研究。

\textbf{會做什麼。} 您（或您的孩子）對著電腦麥克風，念八個英文單字（\textit{she, see, fee, pea, bee, me, knee, cheat}），每個字念兩次，約需 4 分鐘。螢幕不會顯示分數。之後可能有老師聆聽錄音並進行評分。

\textbf{會保存什麼。} 參與者代碼（不用姓名）、\textit{只有在您勾選下方同意欄時才保存的}錄音，以及由錄音測得的數值。本研究不使用相機或影像。資料保存在由 Kiet Pham 管理的研究用電腦中；錄音只有評分老師與研究團隊會聽，不會以電子郵件或通訊軟體傳送，並於 1 年後刪除。研究報告只呈現彙整結果。

\textbf{風險與好處。} 除了幾分鐘的念字外，沒有已知風險。本研究沒有報酬，也沒有直接好處；研究結果可能有助於改進學習者使用的工具。

\textbf{您的權利。} 參加完全出於自願。您或您的孩子可以隨時停止；您也可以要求我們刪除錄音與資料，無須說明原因，也不會影響您（或孩子）的教學或照護。若要退出，請提供參與者代碼並聯絡上述信箱。

\medskip\hrule\medskip
\textbf{同意事項}（孩子由家長或監護人填寫，並詢問孩子本人是否願意參加）

\begin{itemize}[label=$\square$]
\item 我已閱讀（或聽人讀完）以上說明，並且可以提出問題。
\item 我同意我（或我的孩子）參加上述錄音。
\item 我\textbf{也}同意\textbf{保存錄音}，並由評分老師與研究團隊依上述方式聆聽。\textit{（若未勾選，仍可參加，但老師無法對錄音評分。）}
\end{itemize}

\vspace{6pt}
\noindent 參與者代碼（由老師填寫）：\rule{3.5cm}{0.4pt}\hfill 日期：\rule{4cm}{0.4pt}\\[12pt]
\noindent 成人參與者或家長／監護人姓名：\hrulefill\\[12pt]
\noindent 簽名：\hrulefill\\[12pt]
\noindent 說明本研究的老師姓名縮寫：\rule{3cm}{0.4pt}\\[12pt]
\noindent 孩子的同意（「我願意參加」），若孩子能表達：\rule{4cm}{0.4pt}

\vspace{4pt}\footnotesize 請將本同意書與錄音分開存放：同意書上有姓名，資料只有代碼。
\end{document}
